\documentclass[../../main.tex]{subfiles}
\begin{document}

{\bf [解]}
題意より$f(x), g(x)$およびその一階微分は
\begin{align}
 \begin{cases}
  f(x) = x^3+px+q \\
  g(x)=x^3+ax^2+b
 \end{cases} \label{eq:1}\\
 \begin{cases}
  f'(x) = 3x^2 + p \\
  g'(x) = 3x^2 + 2ax
 \end{cases} \label{eq:2}
\end{align}
である．
題意より$f, g$ が極大・極小値を持つので 二次方程式$f'(x) = 0, \, g'(x) = 0$ は各々 2 実解を持つ．

$g'(x)=0$の解は
\begin{align}
 3x^2 + 2ax = 0 \\
 \therefore
 x = -\dfrac{2a}{3}, 0 \label{eq:3}
\end{align}
である．

同様に$f'(x)=0$の解は$p<0$のもとで二つの実数解をもち，その値は
\begin{align}
 3x^2 + p = 0 \\
\therefore
 x = \pm\sqrt{\dfrac{-p}{3}} \label{eq:4}
\end{align}
である．

$a>0$に注意して$f(x), g(x)$の増減表は以下のようになる．

% --- f(x) の増減表 ---
\begin{table}[htb]
\centering
\caption{$f(x)$ の増減表}
\begin{tabular}{c|ccccc}
$x$ & $\dots$ & $-\sqrt{-\frac{p}{3}}$ & $\dots$ & $\sqrt{-\frac{p}{3}}$ & $\dots$ \\ \hline
$f'(x)$ & $+$ & $0$ & $-$ & $0$ & $+$ \\ \hline
$f(x)$ & $\nearrow$ & 極大 & $\searrow$ & 極小 & $\nearrow$
\end{tabular}
\end{table}

% --- g(x) の増減表 ---
\begin{table}[htb]
\centering
\caption{$g(x)$ の増減表}
\begin{tabular}{c|ccccc}
$x$ & $\dots$ & $-\frac{2a}{3}$ & $\dots$ & $0$ & $\dots$ \\ \hline
$g'(x)$ & $+$ & $0$ & $-$ & $0$ & $+$ \\ \hline
$g(x)$ & $\nearrow$ & 極大 & $\searrow$ & 極小 & $\nearrow$
\end{tabular}
\end{table}

以下，$p, q$を$a, b$で表す．
増減表より$f, g$はそれぞれ$x=-\sqrt{\dfrac{-p}{3}},-\dfrac{2a}{3}$で極大，$x=\sqrt{\dfrac{-p}{3}}, 0$で極小をとるので，題意の条件は
\begin{align}
 \begin{cases}
  f\left(-\sqrt{-\frac{p}{3}}\right) = g\left(-\frac{2}{3}a\right) \\
  f\left(\sqrt{-\frac{p}{3}}\right)  = g\left(0\right) 
 \end{cases} \label{eq:5}
\end{align}
と書き下せる．
ここで\cref{eq:1}より
\begin{align}
 \begin{cases}
   g\left(-\frac{2}{3}a\right) = (-\frac{2}{3}a)^3 + a(-\frac{2}{3}a)^2 + b = \frac{4}{27}a^3 + b \\
   g(0) = b
 \end{cases} \label{eq:6}
\end{align}
および
\begin{align}
\begin{cases}
 f(\alpha) + f(\beta) = 2q \\
 f(\alpha) - f(\beta) = -\dfrac{4p}{3}\sqrt{-\dfrac{p}{3}}
\end{cases} \label{eq:7}
\end{align}
だから，\cref{eq:5}の辺々足し引きして\cref{eq:6,eq:7}を代入すると
\begin{align}
 \begin{cases}
    2q = \frac{4}{27}a^3 + 2b \\
    -2\sqrt{-\frac{p}{3}} (\frac{2}{3}p) = \frac{4a^3}{27}
 \end{cases} \\
\therefore
 \begin{cases} 
    q = \frac{2}{27}a^3 + b \\ 
    p = -\frac{a^2}{3}
 \end{cases}
\end{align}
である．



\newpage

\end{document}
